% pdfTeX protrudes the first and last character of an `\hbox`, not only of a
% line the paragraph breaker made. There is no §-number for this — protrusion is
% pdfTeX's, not Knuth's — so the behaviour below is what `pdftex -ini` does, and
% these probes are the specification.
%
% It matters because every equation tag goes through one. amsmath sets `(1.1)`
% in an `\hbox`, and with microtype loaded pdfTeX gives that box a left margin
% kern for the `(` and a right margin kern for the `)`:
%
%     \hbox(8.2125+2.73749)x19.94624, display
%     .\hbox(8.2125+2.73749)x19.94624
%     ..\kern -1.28114
%     ..(1.1
%     ..\kern 0.0
%     ..)
%     ..\kern -1.28114
%
% Without them the box is 22.50853pt — 2.56228pt too wide, which is 2*1.28114
% exactly. That difference moved a line, which moved a hyphenation, which moved
% a word off page 2 of 2608.17397v1.
%
% `\lpcode`/`\rpcode` are thousandths of the font's quad (`\fontdimen6`), not of
% the character's own width. See docs/history/engine-notes/protrusion.md.
\input plain
\catcode`\@=11
\pdfprotrudechars=1
\lpcode\font`\(=200
\rpcode\font`\)=200
% The equation-tag shape: a box whose first and last characters both protrude.
\setbox0=\hbox{(1.1)}
\message{<tag=\the\wd0>}
% One-sided: only the opening character has a code.
\setbox1=\hbox{(AB}
\message{<left=\the\wd1>}
% Only the closing one.
\setbox2=\hbox{AB)}
\message{<right=\the\wd2>}
% Neither character has a code — no kerns at all.
\setbox3=\hbox{AB}
\message{<none=\the\wd3>}
% A box `to` a fixed size still gets the kerns; they change the glue set, not
% the width.
\setbox4=\hbox to 50pt{(1.1)}
\message{<to=\the\wd4>}
% Turning protrusion off leaves the box alone.
\pdfprotrudechars=0
\setbox5=\hbox{(1.1)}
\message{<off=\the\wd5>}
\pdfprotrudechars=1
% A nested box: the outer box's edge character is inside the inner one.
\setbox6=\hbox{\hbox{(1.1)}}
\message{<nested=\the\wd6>}
% A leading glue node means there is no protruding character on the left.
\setbox7=\hbox{\kern3pt(1.1)}
\message{<kernfirst=\the\wd7>}
\end
